% TOP_PROBLEM: {"id": 3099, "title": "Cube-Simplex conjecture", "category_id": 6, "category": {"id": 6, "name": "geometry", "display_name": "Geometry", "description": "Euclidean and non-Euclidean geometry, geometric structures.", "slug": "geometry", "order_index": 6, "created_at": "2026-07-31T15:26:25.671Z"}, "status": "open", "problem_number": "OPG-778", "set_id": 12}
% FIRST_POSTED: 2026-10-01
% ATTEMPT_STATUS: unresolved
% COMPLETION: 5
% MODEL: GPT 6 Astra Ultra
% UnsolvedMath ID 3099; source catalogue code OPG-778
% !TeX program = lualatex
% Research attempt, 2026-10-01. Canonical statement preserved verbatim.
\documentclass[11pt,a4paper]{article}
\usepackage[margin=25mm]{geometry}
\usepackage{fontspec}
\setmainfont{lmroman10-regular.otf}[BoldFont=lmroman10-bold.otf,
 ItalicFont=lmroman10-italic.otf,BoldItalicFont=lmroman10-bolditalic.otf]
\usepackage{amsmath,amssymb,mathtools}
\usepackage{unicode-math}
\setmathfont{latinmodern-math.otf}
\AtBeginDocument{\let\setminus\smallsetminus}
\usepackage{xurl,hyperref}
\hypersetup{colorlinks=true,urlcolor=blue}
\setcounter{secnumdepth}{0}
\setlength{\parindent}{0pt}
\setlength{\parskip}{5pt}
\setlength{\emergencystretch}{3em}
\begin{document}
\section*{Cube-Simplex conjecture}\label{3099}
{\small UnsolvedMath ID 3099 · OPG-778 · Geometry · OpenGarden}
\subsection{Definitions and mathematical statement}
A convex polytope is the convex hull of finitely many points in a
finite-dimensional real affine space. Its dimension is the dimension of
its affine span. A face is the set on which an affine linear functional
attains its maximum on the polytope; the whole polytope is also allowed
as an improper face.

A $k$-simplex is the convex hull of $k+1$ affinely independent points.
A $k$-dimensional face is combinatorially a cube when its face lattice,
ordered by inclusion, is isomorphic to the face lattice of $[0,1]^k$.
This condition fixes incidences, not Euclidean lengths, angles, or a
particular affine realization.

The cube--simplex conjecture is the quantified statement
For every integer $k\ge1$, there is an integer $d\ge k$ such that
 every convex polytope $P$ of dimension at least $d$ has a $k$-face $F$
 satisfying
\[
       F\text{ is a simplex or a combinatorial cube}.
\]
The dimension threshold may depend on $k$ but must be uniform over all
polytopes and all their numbers of vertices. Either alternative may occur;
it need not be the same alternative for all polytopes.

The target is a face of the original polytope. A simplex or cube found
as an arbitrary subset, a section, a projection or a subdivision cell
does not meet the requirement. There is no simplicity or simpliciality
assumption on $P$. The assertion can also be expressed by asking whether
the least such threshold $d$ is finite for every prescribed $k$.
\subsection{Short English statement}
Must every sufficiently high-dimensional convex polytope have a face of a prescribed dimension that is a simplex or combinatorially a cube?
\subsection{Research attempt}
\paragraph{Literature check and scope.}
The catalogue entry [S2] predates the recent preprint by De Loera, Fang, Guo,
Lu and Zheng [A1], whose abstract reports finiteness of the corresponding
threshold for simple polytopes. That is a reported theorem of those authors;
it is not independently proved here and does not establish the unrestricted
conjecture. This attempt finds the exact threshold within the smaller class
of Cartesian products of simplices, illustrating both a positive mechanism
and a sharp obstruction within that class. Sources were checked 1 October 2026.

\paragraph{Faces of a product.}
Let $P=\prod_{i=1}^r\Delta_{d_i}$ with all $d_i\ge1$. A linear functional on
the product is a sum of factor functionals. Its maximum is attained exactly
on the product of their maximizer faces. Thus every nonempty face is
$\prod_i\Delta_{a_i}$, where $0\le a_i\le d_i$, and its dimension is
$\sum_i a_i$. Conversely every such product is exposed by choosing exposing
functionals independently in the factors. Zero-dimensional factors are points
and do not change its combinatorial type.

\paragraph{A sharp restricted threshold.}
Fix $k\ge2$ and suppose
\[
                    \dim P\ge(k-1)^2+1.
\]
If some $d_i\ge k$, take a $k$-face in that simplex and vertices in all other
factors. This gives a simplex $k$-face of $P$. Otherwise every $d_i\le k-1$.
The dimension inequality then forces $r\ge k$. Choose an edge in each of
$k$ factors and vertices in all the rest. The resulting face is a product
of $k$ segments, hence a combinatorial $k$-cube. This proves the conjectural
conclusion for this class with threshold $(k-1)^2+1$.

To see sharpness, take $P_0=(\Delta_{k-1})^{k-1}$, of dimension $(k-1)^2$.
Any $k$-face would have factor dimensions $a_i$ summing to $k$, each at most
$k-1$. At least two are positive. Its number of vertices is
$\prod_i(a_i+1)>1+\sum_i a_i=k+1$, because the product expansion has at least
one positive cross term. Hence it is not a simplex.
It is not a cube either. If every positive $a_i$ were one, their sum would
be at most $k-1$. Thus some $a_i\ge2$, and a triangular $2$-face in that
factor, with vertices in all others, gives a triangular $2$-face of the
candidate. Every $2$-face of a cube is a quadrilateral, a combinatorial
contradiction. For $k=2$ the example simply has dimension one and no $2$-face;
for $k\ge3$ the argument excludes every possible $k$-face explicitly.

\paragraph{Concrete checks.}
For $k=3$, the product of two triangles has dimension four. Its $3$-faces
are triangular prisms, with six vertices and triangular faces, so none is
a tetrahedron or a cube. In dimension five any product of simplices either
has a factor of dimension at least three or has at least three nontrivial
factors, giving the desired face. The accompanying check enumerates the
factor-dimension possibilities in the obstruction products for $3\le k\le6$;
it verifies the vertex-count and triangular-face obstructions separately.
The general proof is the dimension and product-face argument, not this finite
enumeration.

\paragraph{First remaining gap.}
An arbitrary convex polytope need not have any decomposition into Cartesian
factors, and high dimension does not itself supply independent edge choices
whose product is a face. Nor is an embedded cube subgraph sufficient: the
requirement concerns an exposed face and its entire face lattice. The proof
therefore cannot replace the product hypothesis by arbitrary convexity.
The stronger simple-polytope result reported in [A1] also cannot be extended
to nonsimple polytopes merely by perturbing vertices, because perturbation
can destroy the particular face incidences being sought.

\paragraph{Reproducibility.}
The finite calculations above are implemented in
\texttt{scripts/check\_attempt\_batch\_2026\_10\_01.py}; run it with Python 3
from the repository. It uses exact integer arithmetic and no external packages.

\paragraph{Final status.}
\textbf{UNRESOLVED.} The exact threshold for products of simplices is proved;
the full all-polytopes assertion is not. The restricted threshold is not
claimed to be historically new. Completion estimate: 5\%.

\subsection{Sources}
\begin{itemize}
\item[S1] ULAM.AI and the UnsolvedMath Contributors, supplied problems.json, record 3099 (OPG-778), OpenGarden collection. Catalogue identity and source transcription; CC BY 4.0.
\url{https://huggingface.co/datasets/ulamai/UnsolvedMath}.
\item[S2] Open Problem Garden, Cube-Simplex conjecture. Original problem and discussion; the mathematical formulation below is expanded from the supplied source record.
\url{http://www.openproblemgarden.org/op/cube_simplex_conjecture}.
\item[A1] J. A. De Loera, E. X. Fang, S. Guo, J. Lu and H. Zheng, \emph{On Unavoidable Faces of High-Dimensional Polytopes}, arXiv:2609.00397v1, submitted 31 August 2026. Abstract and manuscript consulted for current context; its general theorem is not reproved here. \url{https://arxiv.org/abs/2609.00397}.
\end{itemize}
\end{document}
